\subsection{Small-object coverage and placement}
\label{sec:existing-reconstruction}
Across the first three held-out rooms, Inventory correctly placed \textbf{30 of 174}
small objects by center and size, compared with \textbf{2 of 174} for original
AHa-3D and 25 for All objects (Table~\ref{tab:small-placement}). With box
overlap as the criterion, the count rose from two to 20 to 21.

\begin{table}[!htbp]
\centering\small
\caption{Small objects correctly placed out of 174 in three held-out rooms.
All objects expands the baseline scope. Inventory adds detector evidence.
The last two rows use the same generator, triggered by object size or by the agent. Parentheses
give gains over the baseline with scale fitted from cameras.}
\label{tab:small-placement}
\begin{tabular}{lrrrr}
\toprule
& \multicolumn{2}{c}{Fitted scale} & \multicolumn{2}{c}{Fixed scale} \\
\cmidrule(lr){2-3}\cmidrule(lr){4-5}
Variant & Center and size & Box overlap & Center and size & Box overlap \\
\midrule
AHa-3D (baseline) & 2 & 2 & 2 & 1 \\
All objects & 25 (+23) & 20 (+18) & 12 & 7 \\
Inventory & \textbf{30 (+28)} & \textbf{21 (+19)} & 13 & 7 \\
\midrule
Required generation & 24 (+22) & 21 (+19) & 9 & 7 \\
Selective generation & 26 (+24) & 18 (+16) & 10 & 6 \\
\bottomrule
\end{tabular}
\end{table}

Relative to All objects, Inventory's net gain was five placements by center
and size, one by fitted-scale overlap and zero by fixed-scale overlap.
Below 20\,cm, Boxer's inventory loosely matched only 10 of the 138 targets
in these rooms, and the Inventory reconstruction correctly placed nine.

Neither generation policy improved on Inventory's 30 small-object placements
by center and size. Across both policies, 160 generation attempts yielded ten
accepted placements and five generated assets in the saved scenes. Flat meshes,
incorrect dimensions and poor silhouette fits prompted the agent to fall back
to procedural geometry.

\textbf{Six rooms and the checklist.}
We then ran the baseline, All objects, Inventory and Checklist on three
further held-out rooms, the next in a fixed order, and ran Checklist on the
first three. The six rooms hold 857 targets: 264 under 20\,cm and 279 from
20 to under 50\,cm (Table~\ref{tab:outcomes}). Of these 543 targets under
50\,cm, the baseline correctly placed 6, All objects 71, Inventory 101 and
Checklist 142. Checklist left 257 of them without a matched reconstruction,
against 361 for Inventory.

\begin{table}[!htbp]
\centering\small
\caption{Outcome of each of the 857 targets in six held-out rooms under the
center-and-size test with fitted scale. Each cell gives targets correctly
placed / matched but misplaced / left without a matched reconstruction.}
\label{tab:outcomes}
\begin{tabular}{lcccc}
\toprule
Variant & $<$20\,cm (264) & 20--$<$50\,cm (279) & 50--$<$150\,cm (206) & $\ge$150\,cm (108) \\
\midrule
AHa-3D (baseline) & 0 / 0 / 264 & 6 / 6 / 267 & 47 / 5 / 154 & 30 / 4 / 74 \\
All objects & 21 / 26 / 217 & 50 / 55 / 174 & 82 / 23 / 101 & 45 / 8 / 55 \\
Inventory & 36 / 22 / 206 & 65 / 59 / 155 & 93 / 22 / 91 & 43 / 11 / 54 \\
Checklist & \textbf{50} / 65 / 149 & \textbf{92} / 79 / 108 & 92 / 31 / 83 & \textbf{50} / 12 / 46 \\
\bottomrule
\end{tabular}
\end{table}

Table~\ref{tab:paired} compares variants room by room. All objects beat the
baseline in all six rooms, the strongest result six rooms can give. Checklist
gained 71 placements over All objects, 41 of them in the room with a wall of
postcards, and was ahead of Inventory in three rooms and behind in three.
Neither difference is significant. Moving every reconstructed instance by a
random 0.5 to 1.5\,m left Checklist with about four correct placements under
50\,cm, so the gain is not an effect of building more boxes. Checklist runs
averaged 79 minutes and \$13.28, against 59 minutes and \$12.36 for All
objects and 104 minutes and \$19.95 for Inventory.

\begin{table}[!htbp]
\centering\small
\caption{Difference in targets under 50\,cm correctly placed in each
held-out room. $p$: exact two-sided sign-flip test over rooms; with six
rooms it cannot fall below 0.031.}
\label{tab:paired}
\begin{tabular}{lrrrrrrr}
\toprule
Room (targets) & 1 (38) & 2 (258) & 3 (16) & 4 (71) & 5 (32) & 6 (128) & Total \\
\midrule
All objects $-$ baseline & $+8$ & $+15$ & $+7$ & $+7$ & $+8$ & $+20$ & $+65$ ($p=0.031$) \\
Inventory $-$ All objects & $-1$ & $+7$ & $+1$ & $+7$ & $+2$ & $+14$ & $+30$ ($p=0.094$) \\
Checklist $-$ All objects & $+9$ & $+18$ & $-3$ & $+5$ & $+1$ & $+41$ & $+71$ ($p=0.094$) \\
Checklist $-$ Inventory & $+10$ & $+11$ & $-4$ & $-2$ & $-1$ & $+27$ & $+41$ ($p=0.25$) \\
\bottomrule
\end{tabular}
\end{table}

The agent recorded a decision for all 2{,}287 checklist entries: 851 built,
647 duplicates, 727 parts of larger objects and 62 skipped. Checklist matched
more targets under 50\,cm than Inventory, 286 against 182, but half of its matches
failed the center-and-size test, so placement now limits it as much as
coverage. One Checklist run in room 2 ended without external review when
the reviewer's quota ran out; we repeated it, and the excluded run placed 36
targets under 50\,cm against 33 for the repeat.

Recall ignores reconstructed instances that match no target. In the first
three rooms these were 39 of 208 for All objects and 24 of 198 for
Inventory. One author checked all 63 against the video: 22 were present
where reconstructed but not annotated, 27 were real but too far off to
match, five were repeats and nine did not appear.
Checklist built 922 instances in six rooms, and 451 matched no target: 168
are postcards on one wall and 108 ornaments on one Christmas tree, which
CA-1M does not annotate one by one. The rest are unchecked.

\textbf{A revised inventory on development rooms.}
Table~\ref{tab:inventory} scores inventories on their own in five development
rooms, treating each inventory box as a placed object. Boxer's boxes passed
center and size for 10\% of the 216 targets under 20\,cm, and SAM~3 with a
fixed category list for 13\%. Room-specific categories raised this to 61\%.
The complete revised inventory reached 58\%, and 51\% of the 141 small
targets. With every box shifted horizontally by a random 0.5 to 1.5\,m
before the final step, about 1\% passed.

\begin{table}[!htbp]
\centering\small
\caption{Inventories scored without an agent on five development rooms.
Rows beginning with + add one step to the row above. Matched: share of
targets with a candidate match. Placed: share whose inventory box passes
center and size. Entries matched: share of inventory entries that match any
target.}
\label{tab:inventory}
\begin{tabular}{lrrrrrr}
\toprule
& & Entries & \multicolumn{2}{c}{$<$20\,cm (216)} & \multicolumn{2}{c}{20--$<$50\,cm (141)} \\
\cmidrule(lr){4-5}\cmidrule(lr){6-7}
Inventory & Entries & matched & Matched & Placed & Matched & Placed \\
\midrule
Boxer & 334 & 66\% & 16\% & 10\% & 57\% & 38\% \\
SAM~3, fixed categories & 745 & 34\% & 27\% & 13\% & 69\% & 39\% \\
+ windows, room categories & 2{,}123 & 20\% & 82\% & 61\% & 81\% & 50\% \\
+ MoGe-2 depth & 2{,}295 & 19\% & 81\% & 58\% & 79\% & 55\% \\
+ support and shared size & 2{,}295 & 19\% & 81\% & 58\% & 79\% & 51\% \\
\bottomrule
\end{tabular}
\end{table}

Two cautions apply. The revised inventory lists 2{,}295 entries against
745, and the share that match a target fell from 34\% to 19\%: it contains
duplicates and object parts. One refrigerator door covered with magnets and
photographs holds 112 of the 216 smallest targets; without that room, the
gain is from 18\% to 40\%. On the six held-out rooms, scored after the
Checklist settings were fixed, the revised inventory placed 32\% of targets
under 20\,cm and 41\% of small targets, against 6\% and 25\% for Boxer.
